% Draft — merged into main.tex once results are final.
\section{Introduction}

Perturb-seq screens measure the transcriptome-wide response of single cells to
CRISPR-mediated gene perturbations \citep{dixit2016,norman2019,replogle2022}.
Because the number of possible combinations grows quadratically with the number
of targets, only a small fraction of double perturbations can ever be measured:
the Norman et al.\ K562 CRISPRa screen profiled 131 of the
$\binom{105}{2}=5{,}460$ pairs of its 105 targets \citep{norman2019}.
Computational models that predict the response to unmeasured perturbations---and
in particular the non-additive, genetic-interaction (GI) component of double
perturbations---are therefore a central goal.

Knowledge-graph neural networks such as GEARS \citep{roohani2024} and TxPert
\citep{wenkel2026} and transformer-based single-cell foundation models
\citep{cui2024scgpt,theodoris2023} have been proposed for this task. However,
recent benchmarks found that these models often do not outperform deliberately
simple baselines---the mean training response, the sum of single-perturbation
effects, or a low-rank linear model
\citep{ahlmann2025,csendes2025}. Two lessons from this literature shape our
study: (i) every new architecture must be compared against such baselines on
identical splits, and (ii) the contribution of any biological prior must be
isolated by controls that remove it while keeping everything else fixed.

State-space models (SSMs) such as Mamba \citep{gu2023} process sequences in time
and memory linear in their length, whereas self-attention is quadratic
\citep{vaswani2017}. Treating the transcriptome as a sequence of gene tokens,
SSMs could scale to all measured genes, but genes have no natural order and an
SSM has no notion of which genes are functionally related. We ask whether both
limitations can be addressed by conditioning the SSM's discretization step on a
gene--gene graph, and---more importantly---whether doing so yields better
predictions than simple baselines.

We make three contributions. First, we introduce GCP-Mamba, a bidirectional
selective SSM over all measured genes in which the diffusion of the perturbation
over a co-expression graph modulates the step size $\Delta$ of every gene token.
Second, we benchmark it on the five standard GEARS splits of the Norman data
against five simple baselines, an official GEARS re-run on the same splits, and
five ablations that remove the graph, permute it, destroy the gene ordering, or
remove the recurrence. Third, we document two evaluation pitfalls we
encountered: a perturbation encoding restricted to highly variable genes that
silently discards most targets, and a GI metric that rewards a model which
predicts no change at all.

\section{Methods}

\subsection{Data and splits}
We used the Norman et al.\ K562 CRISPRa Perturb-seq data \citep{norman2019} as
distributed with GEARS (91{,}205 cells, 5{,}045 genes, 7{,}353 control cells;
log-normalised expression), comprising single perturbations of 105 genes and 131
double perturbations. For each condition we computed the
pseudobulk response $\boldsymbol{\delta}_p=\bar{\mathbf{x}}_p-\bar{\mathbf{x}}_{\mathrm{ctrl}}\in\mathbb{R}^{5045}$.
We used the GEARS ``simulation'' splits with seeds 1--5 \citep{roohani2024},
generated with the official \texttt{cell-gears} code. Each split holds out
$\sim$25\% of targets; test conditions are stratified into unseen singles and
doubles with zero, one or two targets seen as singles in training
(\texttt{combo\_seen0/1/2}). Validation conditions were used only for early
stopping. All graph, ordering and embedding inputs below were computed from the
7{,}353 unperturbed control cells, so they are identical across splits and
contain no information about any perturbation response.

\subsection{Co-expression graph and perturbation signal}
We z-scored control-cell expression, computed absolute Pearson correlations
between all gene pairs, kept the $k=20$ strongest neighbours per gene and
symmetrised the result, giving a sparse graph $\mathbf{W}$ (111{,}078 edges). For a
perturbation set $P$ with indicator $\mathbf{1}_P\in\{0,1\}^G$, the graph signal is
a truncated diffusion
\begin{equation}
\mathbf{c}(P)=\log\!\Big(1+s\sum_{k=1}^{3}\beta^{k-1}\,\widetilde{\mathbf{W}}^{k}\mathbf{1}_P\Big),
\qquad \widetilde{\mathbf{W}}=\mathbf{D}^{-1/2}\mathbf{W}\mathbf{D}^{-1/2},
\end{equation}
with $\beta=0.5$ and $s=10^{3}$ (the logarithm compresses a range spanning four
orders of magnitude). $c_i(P)$ is large for genes close to a perturbed gene in the
co-expression graph and zero for genes more than three hops away.

\subsection{GCP-Mamba}
\textbf{Tokens.} Each of the $G=5{,}045$ genes is a token. Gene $i$ has a learned
identity vector, a projection of its control-cell embedding $\mathbf{e}_i$ (64
principal-component loadings of control-cell expression) and of its control mean
expression. A perturbation enters in three ways: (i) a learned flag vector added
to the tokens of the targeted genes---all 105 targets are measured genes, so no
target is dropped (Section~\ref{sec:pitfalls}); (ii) a set embedding
$\mathbf{z}_P=\sum_{g\in P}\big(f(\mathbf{e}_g)+\mathbf{t}_g\big)$ added to every
token, where $f$ is a two-layer MLP and $\mathbf{t}_g$ a per-target vector
initialised at zero. Only targets perturbed in training receive gradient for
$\mathbf{t}_g$; a target never seen in training is represented through
$f(\mathbf{e}_g)$ alone, which is what makes zero-shot prediction possible; and (iii)
the graph signal $c_i(P)$, linearly projected and added to token $i$.

\textbf{Graph-conditioned selective scan.} Each layer applies a Mamba-style
selective SSM \citep{gu2023} with diagonal state matrix
$\mathbf{A}=-\mathrm{diag}(1,\dots,N)$, input-dependent $\mathbf{B}_i$ and
$\mathbf{C}_i$, a SiLU output gate and a skip connection. Its step size is
conditioned on the graph:
\begin{equation}
\boldsymbol{\Delta}_i=\mathrm{softplus}\big(\mathbf{W}_{\Delta}\mathbf{u}_i+\mathbf{b}_{\Delta}+\mathbf{w}_c\,c_i(P)\big),
\end{equation}
where $\mathbf{u}_i$ is the layer input and $\mathbf{w}_c\in\mathbb{R}^{d}$ is
learned (initialised at zero, so the model starts as an unconditioned SSM). With
zero-order-hold discretization,
$\mathbf{h}_i=\exp(\boldsymbol{\Delta}_i\mathbf{A})\,\mathbf{h}_{i-1}+\boldsymbol{\Delta}_i\mathbf{B}_i u_i$:
tokens near the perturbation take larger steps, writing more of their input into
the state and forgetting more of the past. Unlike a global, perturbation-
independent modulation, $\mathbf{w}_c c_i(P)$ differs between genes and between
perturbations, and permuting the graph changes the function the model computes.

\textbf{Order and direction.} Genes are placed along the Fiedler vector
\citep{fiedler1973} of the normalised Laplacian of $\mathbf{W}$, so strongly
co-expressed genes are close in the sequence, and each layer sums a forward and a
backward scan.

\textbf{Readout.} The prediction for gene $i$ is
$\hat\delta_i=b_i+\mathbf{w}^{\top}\mathbf{h}_i+(\mathbf{U}\mathbf{h}_i)^{\top}(\mathbf{V}\mathbf{z}_P)$,
with $\mathbf{h}_i$ the final token state and rank-16 matrices $\mathbf{U},\mathbf{V}$.
The bilinear term expresses a gene-by-perturbation interaction and contains the
low-rank linear model of \citet{ahlmann2025} as a special case. The bias $b_i$ is
initialised at the mean training response and $\mathbf{w}$ and $\mathbf{V}$ at zero,
so training starts from the \emph{train-mean} baseline and learns deviations from
it; without this initialisation the model predicted spurious changes in thousands
of genes and generalised worse than the mean (Supplementary Section~S2).

\textbf{Implementation and complexity.} We implemented the scan in PyTorch as an
exact chunked parallel scan (chunk length 16; within-chunk cumulative sums in
log space, sequential recurrence across chunk states). Per-step log-decays are
clamped at $-5$ ($\bar a\geq 0.0067$) so that the within-chunk exponentials stay
finite in single precision; the implementation matches the naive recurrence to
$10^{-5}$ (unit test in the repository). Time and memory are
$\mathcal{O}(BGdN)$ for the scan plus $\mathcal{O}(|\mathbf{W}|)$ for the sparse
diffusion, i.e.\ linear in the number of genes. No $G\times G$ object is formed.

\textbf{Training.} Two layers, $d=32$, $N=4$, dropout 0.1 in $f$. AdamW
\citep{loshchilov2019}, learning rate $10^{-3}$, weight decay $10^{-2}$,
batch size 8, gradient clipping at 1, mean-squared error on all genes, up to 100
epochs with early stopping on validation MSE (patience 15). The model and data
seed equal the split seed. All experiments ran on a 4-core CPU.

\subsection{Ablations}
All ablations share the architecture, inputs (i)--(ii), optimiser and stopping
rule, and differ in one factor: \emph{$\Delta$ only}, graph used only in
$\Delta$ and not added to tokens; \emph{permuted graph}, the gene labels of
$\mathbf{W}$ are randomly permuted (preserving its degree distribution and
spectrum); \emph{random order}, genes are placed in a random order instead of the
Fiedler order; \emph{no graph}, $c(P)\equiv0$ (a plain bidirectional Mamba);
\emph{Graph-MLP}, the SSM blocks are replaced by per-token MLPs so that genes
interact only through the shared perturbation embedding and the graph signal.

\subsection{Baselines}
\emph{No change} predicts $\boldsymbol{\delta}=\mathbf{0}$. \emph{Train mean}
predicts the mean training response. \emph{Additive} predicts
$\boldsymbol{\delta}_A+\boldsymbol{\delta}_B$ for a double using the measured
training singles, substituting the mean single response for any target without a
training single. \emph{Linear} is the low-rank model of Ahlmann-Eltze et al.\
\citep{ahlmann2025}, $\mathbf{Y}\approx\mathbf{G}\mathbf{M}\mathbf{P}^{\top}+\mathbf{b}$,
with gene embeddings $\mathbf{G}$ from the top 10 principal components of the
training responses, ridge penalty $0.1$ and each perturbation represented by the
sum of its targets' embeddings, taken either from $\mathbf{G}$ (\emph{Linear
(PCA)}, as in the original) or from the control-cell embeddings $\mathbf{e}_g$ that
GCP-Mamba receives (\emph{Linear (ctrl emb.)}), columns standardised. \emph{GEARS} was retrained with
the official implementation and default hyperparameters (hidden size 64, 20
epochs) on the same five splits.

\subsection{Metrics and statistics}
Following \citet{roohani2024}, for each test condition we computed the MSE and the
Pearson correlation of predicted versus observed $\boldsymbol{\delta}$ on its top-20
differentially expressed (DE) genes, and the fraction of those genes whose
direction of change is predicted correctly. For doubles whose singles were both
measured, we computed the GI residual
$\boldsymbol{\epsilon}=\boldsymbol{\delta}_{AB}-\boldsymbol{\delta}_A-\boldsymbol{\delta}_B$
(measured singles, used only for scoring) and report the MSE and Pearson
correlation between predicted and observed residuals on the top-20 DE genes.
Metrics were averaged over test conditions within a split and then summarised as
mean\,$\pm$\,s.d.\ over the five splits. Model comparisons use two-sided Wilcoxon
signed-rank tests on per-condition differences pooled across splits, with
Holm correction across the comparisons reported in each table.
